% from ega3-6-fr.tex
\oldpage[III]{152}
inférieurement (cf. (0, 11.5.2.2)); notre assertion résulte de ce que
l'isomorphisme (6.3.4.1) ne dépend pas de la résolution projective
$L''_{\bullet\bullet}$ choisie, en vertu de l'existence d'un homotopisme entre
deux telles résolutions (M, V, 1.2).
\end{namedtheorem}

\begin{proposition}[6.5.13]
Soient $X$, $Y$ deux $S$-préschémas, et supposons vérifiée l'une des conditions
suivantes :
\begin{enumerate}
  \item[(i)] $X$ et $Z=X\times_S Y$ sont localement noethériens et $X$ est plat
    sur $S$.
  \item[(ii)] $S$ et $X$ sont localement noethériens et $Y$ est de type fini
    sur $S$.
\end{enumerate}
Soit $\mathscr P_\bullet$ (resp. $\mathscr Q_\bullet$) un complexe de
$\mathscr O_X$-Modules (resp. de $\mathscr O_Y$-Modules) quasi-cohérents,
limité inférieurement. On suppose en outre que, pour tout $n$,
$\mathscr H_n(\mathscr P_\bullet)$ (resp.
$\mathscr H_n(\mathscr Q_\bullet)$) est un $\mathscr O_X$-Module (resp. un
$\mathscr O_Y$-Module) de type fini. Alors les
$\mathscr{T}\!or_n^S(\mathscr P_\bullet,\mathscr Q_\bullet)$ sont des
$\mathscr O_Z$-Modules cohérents.
\end{proposition}

\begin{proof}
Comme $\mathscr P_\bullet$ et $\mathscr Q_\bullet$ sont limités
inférieurement, la suite spectrale
$''\mathscr E(\mathscr P_\bullet,\mathscr Q_\bullet)$ est birégulière (6.5.4),
et en vertu de (0, 11.1.8), il suffit (puisque dans les deux cas (i), (ii), $Z$
est localement noethérien) de prouver que les termes
$''\mathscr E_{pq}^2$ sont cohérents. L'hypothèse sur les
$\mathscr H_n(\mathscr P_\bullet)$ et
$\mathscr H_n(\mathscr Q_\bullet)$ et l'expression (6.5.4.2) des
$''\mathscr E_{pq}^2$ montrent donc que la proposition est équivalente à son
cas particulier correspondant à $\mathscr P_\bullet$ et
$\mathscr Q_\bullet$ réduits à leurs termes de degré 0, autrement dit à son
\end{proof}

\begin{corollary}[6.5.14]
Supposons vérifiée l'une des conditions (i), (ii) de (6.5.13), et soit
$\mathscr F$ (resp. $\mathscr G$) un $\mathscr O_X$-Module (resp. un
$\mathscr O_Y$-Module) quasi-cohérent de type fini; alors les
$\mathscr{T}\!or_n^S(\mathscr F,\mathscr G)$ sont des $\mathscr O_Z$-Modules
cohérents.
\end{corollary}

\begin{proof}
La question étant locale sur $X$ et $Y$, on peut supposer $S$, $X$, $Y$
affines.

\textnormal{(i)} Sous les hypothèses de (i), $S$, $X$ et $Z$ sont noethériens.
Il existe donc une résolution localement libre $\mathscr L_\bullet$ de
$\mathscr F$ formée de $\mathscr O_X$-Modules de type fini (I, 1.3.7); comme
$X$ est plat sur $S$, il résulte de (6.3.4) que l'on a
\[
  \mathscr{T}\!or_n^S(\mathscr F,\mathscr G)
  =\mathscr H_n(\mathscr L_\bullet\otimes_S\mathscr G);
\]
or, les $\mathscr L_i\otimes_S\mathscr G$ sont des $\mathscr O_Z$-Modules
quasi-cohérents de type fini (I, 9.1.1), donc cohérents. On en conclut que
$\mathscr H_n(\mathscr L_\bullet\otimes_S\mathscr G)$ est cohérent
($0_{\mathrm I}$, 5.3.4).

\textnormal{(ii)} Supposons maintenant vérifiées les conditions (ii). Comme
l'anneau $A(Y)$ est quotient d'une $A(S)$-algèbre de polynômes à un nombre fini
d'indéterminées (I, 6.3.3), $Y$ est un sous-$S$-préschéma fermé d'un
$S$-préschéma affine $Y'$, plat et de type fini sur $S$; $Y'$ étant noethérien
(I, 6.3.7), il existe une résolution localement libre $\mathscr M_\bullet$ de
$j_*(\mathscr G)$ par des $\mathscr O_{Y'}$-Modules de type fini
($j:Y\to Y'$ étant l'injection canonique); en vertu de (6.5.12),
$\mathscr{T}\!or_n^S(\mathscr F,\mathscr G)$ est l'image réciproque par
$1\times j$ du $\mathscr O_{Z'}$-Module
$\mathscr H_n(\mathscr F\otimes_S\mathscr M_\bullet)$ (où
$Z'=X\times_S Y'$); on voit comme dans (i) que les
$\mathscr F\otimes_S\mathscr M_i$ sont des $\mathscr O_{Z'}$-Modules
cohérents, et l'on en tire encore la conclusion par ($0_{\mathrm I}$, 5.3.4).
\end{proof}

\begin{env}[6.5.15]
La théorie développée ci-dessus pour deux complexes $\mathscr P_\bullet$,
$\mathscr Q_\bullet$ de faisceaux quasi-cohérents sur deux $S$-préschémas
$X$, $Y$ se généralise sans peine au cas où l'on considère un nombre fini
quelconque de $S$-préschémas $X^{(i)}$ ($1\leqslant i\leqslant m$) et sur
chaque $X^{(i)}$ un complexe $\mathscr P_\bullet^{(i)}$ de
$\mathscr O_{X^{(i)}}$-Modules quasi-cohérents; si
$Z=X^{(1)}\times_S X^{(2)}\times_S\cdots\times_S X^{(m)}$, on définit ainsi un
$\mathscr O_Z$-Module quasi-cohérent
$\mathscr{T}\!or_q^S(\mathscr P_\bullet^{(1)},\mathscr P_\bullet^{(2)},
\ldots,\mathscr P_\bullet^{(m)})$. Nous laissons au lecteur le soin de
développer la théorie dans ce cas général et nous nous bornerons à écrire,
pour référence ultérieure, le terme $E_2$ de la seconde suite spectrale
(régulière) dont l'aboutissement est
$\mathscr{T}\!or_\bullet^S(\mathscr P_\bullet^{(1)},
\mathscr P_\bullet^{(2)},\ldots,\mathscr P_\bullet^{(m)})$ :

